% !TeX root = ../papers/nc-p3.tex
\section{二次行列鉛筆の核直線と次の分岐中心}
\label{sec:rank-ten-flat-kernel-lines}
\begin{theorem}\label{thm:rank-ten-flat-kernel-lines}
定理\ref{thm:rank-ten-cubic-after-center}の四次元順中心 $p$ と第二ねじれ三次曲線の上に、
その後続直線が一般核直線の極限に一致し、自身の第二曲線が平面二次曲線である順中心 $r$ が存在する。
$p$ の第三曲線は直線または平面二次曲線であり、二次曲線なら第四曲線は直線である。
II型の制限では第三曲線は必ず直線である。
従って反対階数10で曲面射影のない窓のどの四次元順中心からも、
高々三辺で中央支持3の順中心に達する。逆向きでも成立する。
\end{theorem}
\begin{proof}[特殊点での核の極限]
同定理の両標準形で $n=t_0$ とし $t_*=(0,1)$ とする。一般核の極限 $K(t)$ は、
I型では二つの二座標組をそれぞれ $(t_0,t_1)$ が消す二平面、
II型では第四座標が零で残る三座標を $(t_0^2,t_0t_1,t_1^2)$ が消す二平面である。
いずれも全 $t$ で部分束を与える。$E(t):V_3^*\to R_0$ を $x_2(t)$ での順評価とする。
第一直線に $R_0$ を制限した後の $E(t)$ は $[-,t]\otimes A(t)$ であり、
制限 $R_0\to R_0|_{L}$ の核は一次元であった。$t\ne t_*$ で順中心があれば
$E(t)$ は階数2なのでその核は $K(t)$ そのものである。

そうでなければ順側欠損 $f_2=6-d_3\ge2$ は全て $t_*$ にある。
この点の局所一様化変数を $v$ とし、局所自由な基底で $E(v)$ を読む。
制限後の二つの有効な行は $L_1(v),vL_2(v)$ と書け、
$L_1,L_2$ は全点で独立な行で $K=\ker(L_1,L_2)$ となる。
実際両標準形の $u,v(t)$ の方向は全点で独立、二行の一方のみに $n$ が掛かる。
したがって単元な主元で消去し、一次元核側の行 $C(v)$ を戻すと、
三次小行列のイデアルは $v$ と、$K$ の局所基底における $C|_K$ の二成分のイデアルとの積である。
これは $L_1,L_2$ を座標行に取れば直ちに従う。位数が2以上のため $C(0)$ は $K(0)$ を消し、
$E(0)$ の階数は下界により2。従ってここでも $\ker E(0)=K(0)$ である。
こうしていずれの場合にも所望の核直線を持つ順中心 $r$ が得られる。
\end{proof}
\begin{proof}[二つの列の平面切断]
I型では $K(t)$ の点を
$z=(-t_1X,t_0X,-t_1Y,t_0Y)$ とでき、$y$ を独立に動かすと
\[
              A(y)z=[t,y]\big(y_0 X u+Y(y_0v_0+y_1v_1)\big).
\]
従って $K(t)$ の一般点のどの後続も $H=\{v_1=0\}$ に入る。
II型では $K(t)$ の第四座標は零で、$A(y)z$ は恒等零でない $u$ の倍となるから
全ての後続が $H=\{u=0\}$ に入る。特に $r$ の第二曲線は $H$ にある。
第二次数は2または3で、3ならねじれ三次となるから平面二次曲線である。

一つの平面内にある異なる動く成分の次数和は3以下である。
系\ref{cor:rank-nine-nonlinear-corridors}の末尾の単一成分の証明で、
その平面内の異なる曲線の式の積を使えばよい。総次数4以上なら順・逆の全四次小行列が
その積で割れ、全て恒等零または定数比例となる。前者は曲面射影、
後者は同じ平面の法線に始まり終わる連続純関係を与える。
$I$ 型で $p$ の第三曲線が非平面的なら、同定理の射影式より第四曲線は $H$ の中の
次数2以上の別の曲線であり、$r$ の二次曲線との次数和に反する。
従って第三曲線は滑らかな二次曲面の平面切断にあり、直線または平面二次である。
二次なら二次曲面の両射影が非定数なので同じ式から第四曲線は $H$ 内で動き、
次数和により直線となる。

II型では $p$ の第三曲線も $r$ の後続直線も位置3の第四座標の零平面にある。
後者と前者の正規化の交叉長は高々1である。
実際 $r$ の後続直線の上では逆評価が全点で階数3で、その前点は定数 $r$ である。
この直線に入る任意の微小因子で $p$ の第二点はよって $r$ に等しく、
第二曲線のパラメータ表示が閉埋込みなのでその因子長は高々1である。
二つの列の成分は異なるから、平面の直線を制限した因子の次数は第三曲線の次数に等しい。
従ってこの次数は1である。

最後に $p$ の第二曲線が最初から二次なら定理\ref{thm:rank-ten-conic-after-center}で
一辺先に三次元順中心がある。ねじれ三次なら $r$ が三次元であるか、
四次元で自身の第二曲線が二次であり、その後続直線に三次元順中心を持つ。
二項中心の排除と二直線列の排除でこれらの場合が尽きる。
\end{proof}
